\documentclass[a4paper]{article}

%----------------------------------------------------------------------------------------
% Dokumentensprache
%----------------------------------------------------------------------------------------
\usepackage[ngerman]{babel} %%change from ngerman to english if preferred
\usepackage[utf8x]{inputenc}
\usepackage[T1]{fontenc}

%----------------------------------------------------------------------------------------
% Seiteneinrichtung
%----------------------------------------------------------------------------------------
\usepackage[a4paper,top=3cm,bottom=2cm,left=3cm,right=3cm,marginparwidth=1.75cm]{geometry}

%----------------------------------------------------------------------------------------
% Packages
%----------------------------------------------------------------------------------------
\usepackage{amsmath}
\usepackage{graphicx}
\usepackage[colorinlistoftodos]{todonotes}
\usepackage[toc,page]{appendix}
\usepackage[hidelinks]{hyperref}
\usepackage{url}
\usepackage{wrapfig}
\usepackage{eurosym}
\usepackage{xcolor}
\usepackage[firstpage=true]{background}
\usepackage{mwe}
\usepackage{lipsum}

%----------------------------------------------------------------------------------------
% FH Kiel Footer
%----------------------------------------------------------------------------------------
\backgroundsetup{placement=bottom,nodeanchor=north west, scale=1,contents={
\includegraphics[width=0.99\paperwidth]{Logo_FH_Kiel_2Zeiler}} 
}

%----------------------------------------------------------------------------------------
% Dokument
%----------------------------------------------------------------------------------------
\begin{document}
\input{titlepage} %% Titelseite zu bearbeiten in titlepage.tex


\newpage
\section{Grundlagen}
Referenzen auf Abbildungen werden mit Abbildung \ref{fig:fh_kiel} gebildet.\\
Referenzen auf Quellen werden mittels \cite{lit:datasheet} erstellt. Die Quellen sind in der "quellenangaben.bib" Datei hinzuzufügen.\\
Empfohlen werden zur Benennung der Inhalte und leichterer Referenzierung folgende Formate:\\
Abbildungen = fig:abbildungsname\\
Tabellen = tab:tabellenname\\
Quellen = lit:literaturname\\
Anhänge = app:anhangsname\vspace*{3\baselineskip}
\lipsum[1] \cite{lit:datasheet}
%% Normale Grafiken
\begin{figure}[h!]
\centering
\includegraphics[width=0.5\textwidth]{Fachhochschule_Kiel}
\caption{\label{fig:fh_kiel}Beispielgrafik}
\end{figure}
\subsection{Grundlagen 1}
%% Textumflossene Grafiken
\par
\begin{wrapfigure}[4]{r}[0cm]{3cm}
\includegraphics[width=0.2\textwidth]{Fachhochschule_Kiel}
\label{fig:fh_kiel_floating}
\end{wrapfigure}
\lipsum[1] \cite{lit:schutzarten}
\subsection{Grundlagen 2}
\lipsum[1]


\newpage
\section{Durchführung}
\lipsum[1-2]
%% Tabelle
\begin{table}
\centering
\begin{tabular}{|l|l|l|}
   \hline
   Gerät & Typ  & Farbe\\
   \hline
   Fernseher & Sony & Mittelblau\\
   \hline
   \end{tabular}
 \caption{Bespieltabelle}
\end{table}


\newpage
\section{Auswertung}
\lipsum[1-5]
\newpage
\section{Fazit}
\lipsum[1-2]
\appendix
\input{appendix.tex} %% Anhänge zu bearbeiten in appendix.tex
\end{document}